% Part: counterfactuals
% Chapter: introduction
% Section: paradoxes-material

\documentclass[../../../include/open-logic-section]{subfiles}

\begin{document}

\olfileid[pt]{cnt}{int}{par}

\olsection{Paradoxos do condicional material}

Um dos primeiros a criticar o uso de $!A \lif !B$ como modo de
simbolizar enunciados em inglês da forma ``se \dots, então \dots'' foi
C.~I. Lewis. Lewis criticava o uso do condicional material nos
\emph{Principia Mathematica} de Whitehead e Russell, que
pronunciavam $\lif$ como ``implica''. Lewis observou, com razão, que, se
$\lif$ significasse ``implica'', então qualquer proposição falsa~$p$ implicaria que
$p$ implica~$q$, pois $p \lif (p \lif q)$ é verdadeiro se~$p$ é falso;
e qualquer proposição verdadeira~$q$ implicaria que $p$ implica~$q$, pois
$q \lif (p \lif q)$ é verdadeiro se $q$~é verdadeiro.

É claro que os lógicos sabem que a \emph{implicação}, isto é, a
consequência lógica, não é um conectivo, mas uma relação entre !!{formula}s ou
enunciados. Para evitar confusão, basta não ler $\lif$ como ``implica''.
\footnote{A leitura de ``$\lif$'' como ``implica'' ainda é muito
  praticada por matemáticos e cientistas da computação, embora
  filósofos procurem evitar as confusões destacadas por Lewis
  pronunciando-o como ``somente se''.} Desde que não façamos isso, a preocupação
específica de Lewis simplesmente não surge: $p$ não ``implica'' $q$,
mesmo que interpretemos $p$ como uma sentença falsa do inglês. Para
determinar se $p \Entails q$, precisamos considerar \emph{todas} as
!!{valuation}s, e $p \Entails/ q$ mesmo quando usamos $p$ para simbolizar
uma sentença que, de fato, é falsa.

Mas ainda há algo estranho em enunciados da forma ``se \dots,
então \dots'', como o de Lewis:
\begin{quote}
Se a Lua é feita de queijo verde, então $2+2=4$.
\end{quote}
e nas inferências
\begin{quote}
  A Lua não é feita de queijo verde. Portanto, se a Lua é feita
  de queijo verde, então $2+2=4$.

  $2+2 = 4$. Portanto, se a Lua é feita
  de queijo verde, então $2+2=4$.
\end{quote}
No entanto, se ``se \dots, então \dots'' fosse apenas $\lif$, a sentença seria
verdadeira sem qualquer problema, e as inferências seriam válidas sem qualquer problema.

Outro exemplo diz respeito à tautologia $(!A \lif !B) \lor (!B \lif
!A)$. Isso sugeriria que, se tomarmos ao acaso duas sentenças
indicativas~$S$ e $T$ de um jornal, a sentença ``Se
$S$, então $T$; ou, se $T$, então~$S$'' deveria ser verdadeira.

\end{document}
